\section*{Problem 1}

给定线性映射 $T : \mathbb{R}^3 \to \mathbb{R}^3$，其在标准基下对应的矩阵为
\[
A =
\begin{bmatrix}
1 & 2 & 3 \\
2 & 4 & 6 \\
1 & 2 & 3
\end{bmatrix}.
\]
请计算矩阵 $A$ 的秩 $\rank \left(A\right)$，并依此求出 $T$ 的像空间维数 $\dim \im \left(\mathcal{T}\right)$ 与核空间维数 $\dim \ker\left(\mathcal{T}\right)$，最后说明其计算结果是否满足秩零化度定理.

\begin{proof}
对 $A$ 施 Gauss - Jordan 消元法有
\[
A = \begin{bmatrix}
    1 & 2 & 3 \\
    2 & 4 & 6 \\
    1 & 2 & 3 \\
\end{bmatrix}
\implies
R = \begin{bmatrix}
    1 & 2 & 3 \\
    0 & 0 & 0 \\
    0 & 0 & 0 \\
\end{bmatrix}
\implies
\boxed{\begin{cases}
    \rank \left(A\right) = 1, \\
    \dim \im \left(\mathcal{T}\right) = \dim C \left(A\right) = 1, \\
    \dim \ker \left(\mathcal{T}\right) = \dim N \left(A\right) = 2. \\
\end{cases}}
\]
注意到 $\dim \im \left(\mathcal{T}\right) + \dim \ker\left(\mathcal{T}\right) = 1 + 2 = 3 = \dim \mathbb{R}^3$，因此 $\boxed{\text{满足秩零化度定理}}$.

\end{proof}

\section*{Problem 2}

设线性变换 $T : \mathbb{R}^2 \to \mathbb{R}^2$ 定义为
\[
T\begin{pmatrix} x_1 \\ x_2 \end{pmatrix}
=
\begin{pmatrix}
2x_1 + x_2 \\
x_1 + 2x_2
\end{pmatrix}.
\]
\begin{enumerate}
\item 求 $T$ 在标准基 $\bvec{e}_1 = \begin{pmatrix}1\\0\end{pmatrix},\;
\bvec{e}_2 = \begin{pmatrix}0\\1\end{pmatrix}$ 下对应的矩阵 $A$；
\item 给定新基（已固定顺序）
\[
\tilde{\bvec{e}}_1 = \begin{pmatrix}1\\1\end{pmatrix}, \quad
\tilde{\bvec{e}}_2 = \begin{pmatrix}-1\\1\end{pmatrix},
\]
求从标准基到新基的过渡矩阵 $P$，并计算 $T$ 在该新基下对应的矩阵 $B$.
\end{enumerate}

\begin{proof}
依据线性变换性质进行代换
\[
\begin{aligned}
    T \begin{pmatrix}
        x_1 \\
        x_2 \\
    \end{pmatrix} &= T \left(\alpha_1 \bvec{e}_1 + \alpha_2 \bvec{e}_2\right), \\
    &= \alpha_1 T \left(\bvec{e}_1\right) + \alpha_2 T \left(\bvec{e}_2\right) = \begin{bmatrix}
        T \left(\bvec{e}_1\right) & T \left(\bvec{e}_2\right)
    \end{bmatrix} \begin{pmatrix}
        \alpha_1 \\
        \alpha_2 \\
    \end{pmatrix}, \\
    &= \begin{bmatrix}
        \bvec{e}_1 & \bvec{e}_2 \\
    \end{bmatrix} \underbrace{\begin{bmatrix}
        2 & 1 \\
        1 & 2 \\
    \end{bmatrix}}_{A} \begin{pmatrix}
        \alpha_1 \\
        \alpha_2 \\
    \end{pmatrix}.
\end{aligned}
\]
从而
\[
A = \boxed{\begin{bmatrix}
    2 & 1 \\
    1 & 2 \\
\end{bmatrix}}.
\]
按新基向量由旧基坐标组成矩阵的约定，过渡矩阵 $P$ 满足
\[
\begin{bmatrix}
    \tilde{\bvec{e}_1} & \tilde{\bvec{e}_2} \\
\end{bmatrix}
= \begin{bmatrix}
    \bvec{e}_1 & \bvec{e}_2 \\
\end{bmatrix} P
\implies
P = \boxed{\begin{bmatrix}
    1 & -1 \\
    1 & 1 \\
\end{bmatrix}}, P^{-1} = \begin{bmatrix}
    1/2 & 1/2 \\
    -1/2 & 1/2 \\
\end{bmatrix}.
\]
其中 $P$ 将新基坐标变为标准坐标，而 $P^{-1}$ 将标准坐标变为新基坐标。
在两组基下的坐标满足
\[
\begin{aligned}
    \begin{bmatrix}
        \tilde{\bvec{e}_1} & \tilde{\bvec{e}_2} \\
    \end{bmatrix} \begin{pmatrix}
        \tilde{\alpha}_1 \\
        \tilde{\alpha}_2 \\
    \end{pmatrix} = \begin{bmatrix}
        \bvec{e}_1 & \bvec{e}_2 \\
    \end{bmatrix} \begin{pmatrix}
        \alpha_1 \\
        \alpha_2 \\
    \end{pmatrix}
\implies
P \begin{pmatrix}
    \tilde{\alpha}_1 \\
    \tilde{\alpha}_2 \\
\end{pmatrix} = \begin{pmatrix}
    \alpha_1 \\
    \alpha_2 \\
\end{pmatrix}.
\end{aligned}
\]
因此
\[
\begin{aligned}
    \begin{bmatrix}
        \bvec{e}_1 & \bvec{e}_2 \\
    \end{bmatrix} A \begin{pmatrix}
        \alpha_1 \\
        \alpha_2 \\
    \end{pmatrix} = \mathcal{T} \left(\alpha_1 \bvec{e}_1 + \alpha_2 \bvec{e}_2\right) &= \mathcal{T} \left(\bvec{x}\right), \\
    &= \begin{bmatrix}
        \tilde{\bvec{e}_1} & \tilde{\bvec{e}_2} \\
    \end{bmatrix} P^{-1} A \begin{pmatrix}
        \alpha_1 \\
        \alpha_2 \\
    \end{pmatrix} = \begin{bmatrix}
        \tilde{\bvec{e}_1} & \tilde{\bvec{e}_2} \\
    \end{bmatrix} \underbrace{P^{-1} A P}_{B} \begin{pmatrix}
        \tilde{\alpha}_1 \\
        \tilde{\alpha}_2 \\
    \end{pmatrix}.
\end{aligned}
\]
最终得到 $B = \begin{bmatrix}
    3 & 0 \\
    0 & 1 \\
\end{bmatrix}$.


\end{proof}
\section*{Problem 3}

设矩阵
\[
A =
\begin{bmatrix}
4 & -2 \\
1 & 1
\end{bmatrix}.
\]
求矩阵 $A$ 的所有特征值.对于每一个特征值，求其对应的特征子空间的维数（即几何重数），并据此判断矩阵 $A$ 是否可对角化.

\begin{proof}
为求 $A$ 的特征值 $\lambda_i$，解方程
\[
\det \left(\lambda I - A\right) = 0 \iff \det \begin{bmatrix}
    \lambda - 4 & 2 \\
    -1 & \lambda - 1 \\
\end{bmatrix} = \lambda^2 - 5 \lambda + 6 = 0 \iff \begin{cases}
    \lambda_1 = 2, \\
    \lambda_2 = 3. \\
\end{cases}
\]
对于 $\lambda_1 = 2$，解方程
\[
\left(\lambda_1 I - A\right) \bvec{x} = \bvec{0} \iff \begin{bmatrix}
    -2 & 2 \\
    -1 & 1 \\
\end{bmatrix} \bvec{x} = \bvec{0} \implies \bvec{x} = t \begin{pmatrix}
    1 \\
    1 \\
\end{pmatrix}, t \in \mathbb{R}.
\]因此 $\lambda_1$ 的特征子空间为 $\Span \left\{\begin{pmatrix}
    1 \\
    1 \\
\end{pmatrix}\right\}$，其维数为 $1$.同样地，对于 $\lambda_2 = 3$，解方程
\[
\left(\lambda_2 I - A\right) \bvec{x} = \bvec{0} \iff \begin{bmatrix}
    -1 & 2 \\
    -1 & 2 \\
\end{bmatrix} \bvec{x} = \bvec{0} \implies \bvec{x} = t \begin{pmatrix}
    2 \\
    1 \\
\end{pmatrix}, t \in \mathbb{R}.
\]因此 $\lambda_2$ 的特征子空间为 $\Span \left\{\begin{pmatrix}
    2 \\
    1 \\
\end{pmatrix}\right\}$，其维数为 $1$.由于两个特征值的几何重数之和等于矩阵的阶数（即 $1 + 1 = 2$），因此 $\boxed{A \text{ 可对角化}}$.

更直接的方式是给出对角化所用矩阵 $P = \begin{bmatrix}
    1 & 2 \\
    1 & 1 \\
\end{bmatrix}$，从而
\[
AP = P \Lambda \iff P^{-1}AP = \Lambda = \diag \left(2, 3\right).
\]
\end{proof}
